% Supplemental working note. Source parameters and relations reconstructed
% from the Gomez--Venegas manuscript table tab:Mt, at h=1.
\section{An exact ozone computation for $\mathbb Q(1)$}
Let $B=\mathbb Q(1)$ and use the PBW ordering
$x_1^a x_2^b y_1^c y_2^d$. In the localised presentation of
Table~\texttt{tab:Mt} we have $n=x_1$, $t=x_1^{-1}x_2$,
\[
 y\,x_1=x_1\begin{pmatrix}0&1\\-1&0\end{pmatrix}y,
 \qquad y\,t=\begin{pmatrix}-1&1-t\\1+t&1\end{pmatrix}y.
\]
These formulas, together with $x_2x_1=-x_1x_2$ and
$y_2y_1=-y_1y_2$, determine all four mixed relations.
The following element of bidegree $(1,2)$ is central:
\[
 z=x_1(y_1^2+y_2^2-2y_1y_2)+x_2(y_1^2-y_2^2).
\]
Its commutators with $x_1,x_2,y_1,y_2$ vanish by direct PBW
reduction; this is independently rechecked by
\texttt{qrv\_ozone\_h1.py} and \texttt{qrv\_center\_low.py}.

\begin{proposition}
For the above member of the $\mathbb Q$ family,
$\operatorname{Oz}(B)\cong\mathbb Z/4$, generated by
$\operatorname{Ad}_{x_1}$.
\end{proposition}
\begin{proof}
Every ozone automorphism preserves the two degree-one planes by
bigradedness. On each $(-1)$-quantum plane it acts by a diagonal or
antidiagonal matrix. The known central elements
\[
 s_x=x_1^2+x_2^2,\quad
 v=x_1^2x_2^2+\frac12x_2^4,\quad
 s_y=y_1^2+y_2^2
\]
force the $x$-matrix to be diagonal with diagonal entries in
$\{\pm1\}$, and the $y$-matrix to be a signed diagonal or signed
antidiagonal matrix. Testing the four mixed relations leaves exactly
8 candidates, which are enumerated exactly by the certificate.
The additional central element $z$ reduces these candidates to exactly
4. Finally $x_1$ is normal, and conjugation by $x_1$ has matrices
\[
 X=\operatorname{diag}(1,-1),\qquad
 Y=\begin{pmatrix}0&-1\\1&0\end{pmatrix}.
\]
The four powers of this conjugation are distinct and comprise all four
remaining candidates. Therefore the upper and lower bounds coincide.
\end{proof}

\section{A rigorous upper bound for $\mathbb V(1)$}
For $B=\mathbb V(1)$, $x_2$ is central and the central element
$v=x_1^2(x_1-x_2)^2$ is given in the original manuscript. An ozone
automorphism is bigraded. Fixing $x_2$ and $v$ forces its action on
$\kk x_1\oplus\kk x_2$ to be either
\[
 I\qquad\text{or}\qquad
 \begin{pmatrix}-1&1\\0&1\end{pmatrix}
\]
under the row-action convention $g(x_i)=\sum_j X_{ij}x_j$.
Preserving the quantum $(-1)$-plane on the $y$'s and fixing
$y_1^2+y_2^2$ forces a signed permutation matrix. The six quadratic
relations leave precisely four matrix pairs, implying
\[
 |\operatorname{Oz}(\mathbb V(1))|\le 4.
\]
They fix the central elements calculated in several low bidegrees, but
this does not establish pointwise fixation of the full center.
Consequently no equality is asserted here.

\paragraph{Remaining tasks.} Compute a complete center or exhibit an
appropriate normal conjugating element for the order-four candidate
in $\mathbb V(1)$; then extend both families to arbitrary root-of-unity
$h$. The degree-two-normal-element presentation of $\mathbb R$
remains separate and no formula is inferred for it.
